\subsection{Sprint 2}

Após a apresentação remota feita pelos \ac{ages} IV para as \textit{stakeholders}, percebemos que estávamos no caminho certo. Durante a \textit{planning}, minha \textit{squad} ficou responsável pela página de Módulos.

Na mesma aula, juntamente com o \ac{ages} I Andrei Rech, desenvolvemos o componente principal da tela, o \textit{card}, que futuramente seria utilizado em outras telas. Ao fim da aula, distribuímos as tarefas para agilizar o processo, ficando eu responsável por toda a parte de ajustar a \textit{Modal} que seria utilizada para os \ac{crud}.

Estudei para entender como customizar a \textit{Modal} adequadamente. Comecei criando um esqueleto com alguns componentes que já tinham sido desenvolvidos pela equipe na sprint passada. Ajustei detalhes importantes, como a lógica para alterar a imagem do Módulo. Apesar de ter demorado um pouco, consegui fazer com que os dados alterados na \textit{Modal} refletissem corretamente na tela principal.

No início da última semana, minha \textit{squad} conseguiu finalizar as tarefas atribuídas. Decidimos então começar a desenvolver a tela de Cursos, que, no fim, já estava quase pronta, pois era muito semelhante à de Módulos.

No final desta sprint, as enchentes que assolaram grande parte do estado do Rio Grande do Sul afetaram nossa entrega. O prazo foi postergado até que fosse estabelecido um retorno, inicialmente remoto, para que pudéssemos apresentar.

No dia da apresentação, apenas uma das \textit{stakeholders} compareceu, o que desanimou um pouco a equipe. Após tanto esforço e tempo dedicados, a ausência das demais \textit{stakeholders} fez parecer que não estavam motivadas para ver o progresso do projeto.